Este informe analiza la brecha entre la complejidad teórica y el rendimiento práctico de algoritmos de ordenamiento y multiplicación de matrices mediante un estudio experimental en C++. Se implementaron \texttt{Merge Sort}, \texttt{Quick Sort},\texttt{std::sort}, así como los algoritmos \textit{Naive} y \textit{Strassen} para matrices. Las pruebas, realizadas con datos de distintos tamaños, órdenes y dominios, midieron tiempo y memoria. Los resultados muestran que \texttt{std::sort} supera en la mayoría de casos, aunque algoritmos con peor complejidad pueden ser competitivos en escenarios específicos. En multiplicación, \textit{Strassen} es más rápido que el método clásico para $N \geq 256$, a costa de mayor uso de memoria. En conclusión, la elección del algoritmo óptimo depende no solo de la complejidad asintótica, sino también de los datos y las restricciones del entorno.